% !TeX program = xelatex
% !TeX encoding = UTF-8
\documentclass[12pt,openany,oneside]{book}

\usepackage[a4paper,top=25mm,bottom=25mm,left=27mm,right=27mm,headheight=16pt]{geometry}
\usepackage{amsmath,amssymb,amsfonts,mathtools,bm}
\usepackage{amsthm}
\usepackage{booktabs,longtable,array,multirow,tabularx}
\usepackage{graphicx}
\usepackage{xcolor}
\usepackage{tikz}
\usetikzlibrary{arrows.meta,positioning,calc,fit,shapes.geometric,decorations.pathreplacing,matrix}
\usepackage[most]{tcolorbox}
\usepackage{enumitem}
\usepackage{listings}
\usepackage{caption}
\usepackage{subcaption}
\usepackage{float}
\usepackage{fancyhdr}
\usepackage{titlesec}
\usepackage{setspace}
\usepackage{microtype}
\usepackage{hyperref}
\usepackage{xepersian}

\settextfont[Scale=1.08]{Amiri}
\setlatintextfont{DejaVu Serif}
\setmonofont{DejaVu Sans Mono}
\setdigitfont{Amiri}

\definecolor{DeepBlue}{HTML}{163A5F}
\definecolor{Accent}{HTML}{B34B35}
\definecolor{SoftBlue}{HTML}{EAF2F8}
\definecolor{SoftGold}{HTML}{FFF6DF}
\definecolor{SoftGreen}{HTML}{EAF7EF}
\definecolor{SoftRed}{HTML}{FCEDEC}
\definecolor{SoftGray}{HTML}{F4F5F6}
\definecolor{CodeGray}{HTML}{F7F7F7}

\hypersetup{
  colorlinks=true,
  linkcolor=DeepBlue,
  citecolor=Accent,
  urlcolor=Accent,
  pdftitle={پردازش تصویر در دهه آینده: از عمق ریاضیات تا مهندس حرفه‌ای},
  pdfauthor={جزوه درس کارشناسی ارشد}
}

\setstretch{1.22}
\setlength{\parskip}{0.35em}
\setlength{\parindent}{1.3em}
\setlist[itemize]{itemsep=0.25em,topsep=0.35em}
\setlist[enumerate]{itemsep=0.25em,topsep=0.35em}

\pagestyle{fancy}
\fancyhf{}
\fancyhead[LE,RO]{\small\thepage}
\fancyhead[RE]{\small پردازش تصویر در دهه آینده}
\fancyhead[LO]{\small\nouppercase{\leftmark}}
\renewcommand{\headrulewidth}{0.4pt}

\titleformat{\chapter}[display]
  {\normalfont\huge\bfseries\color{DeepBlue}}
  {\filleft\Large فصل \thechapter}
  {8pt}
  {\filright}
\titleformat{\section}{\Large\bfseries\color{DeepBlue}}{\thesection}{0.7em}{}
\titleformat{\subsection}{\large\bfseries\color{Accent}}{\thesubsection}{0.7em}{}

\newtheorem{definition}{تعریف}[chapter]
\newtheorem{theorem}{قضیه}[chapter]
\newtheorem{proposition}{گزاره}[chapter]
\newtheorem{lemma}{لم}[chapter]
\newtheorem{example}{مثال}[chapter]
\newtheorem{remark}{نکته}[chapter]

\newtcolorbox{learningbox}{
  colback=SoftBlue,colframe=DeepBlue,title=اهداف یادگیری,
  fonttitle=\bfseries,breakable,enhanced}
\newtcolorbox{keybox}[1][]{
  colback=SoftGold,colframe=Accent,title={#1},
  fonttitle=\bfseries,breakable,enhanced}
\newtcolorbox{futurebox}{
  colback=SoftGreen,colframe=green!50!black,title=پل به دهه آینده,
  fonttitle=\bfseries,breakable,enhanced}
\newtcolorbox{warningbox}{
  colback=SoftRed,colframe=red!65!black,title=خطای رایج,
  fonttitle=\bfseries,breakable,enhanced}
\newtcolorbox{workshopbox}[1][]{
  colback=SoftGray,colframe=black!55,title={کارگاه خانگی پردازش تصویر: #1},
  fonttitle=\bfseries,breakable,enhanced}
\newtcolorbox{videobox}{
  colback=white,colframe=DeepBlue,title=تقسیم پیشنهادی ویدئوهای ۱۵ دقیقه‌ای,
  fonttitle=\bfseries,breakable,enhanced}

\newcommand{\R}{\mathbb{R}}
\newcommand{\C}{\mathbb{C}}
\newcommand{\N}{\mathbb{N}}
\newcommand{\E}{\mathbb{E}}
\newcommand{\Prob}{\mathbb{P}}
\newcommand{\norm}[1]{\left\lVert #1\right\rVert}
\newcommand{\abs}[1]{\left\lvert #1\right\rvert}
\newcommand{\inner}[2]{\left\langle #1,#2\right\rangle}
\newcommand{\trans}{^{\mathsf T}}
\newcommand{\lrterm}[1]{\lr{#1}}
\newcommand{\vect}[1]{\bm{#1}}
\newcommand{\op}[1]{\mathcal{#1}}

\lstset{
  language=Python,
  basicstyle=\small\ttfamily,
  backgroundcolor=\color{CodeGray},
  frame=single,
  rulecolor=\color{black!25},
  numbers=left,
  numberstyle=\tiny,
  breaklines=true,
  showstringspaces=false,
  keywordstyle=\bfseries,
  commentstyle=\itshape,
  tabsize=4
}

\begin{document}

% -----------------------------------------------------------------------------
% Title page
% -----------------------------------------------------------------------------
\begin{titlepage}
\centering
\vspace*{1.2cm}
{\color{DeepBlue}\rule{0.88\textwidth}{1.1pt}}\\[1.2cm]
{\Huge\bfseries پردازش تصویر در دهه آینده\\[0.35cm]}
{\LARGE\bfseries از عمق ریاضیات تا مهندس حرفه‌ای\\[0.9cm]}
{\Large کتاب ـ جزوه درس کارشناسی ارشد و آزمایشگاه پژوهش‌محور\\[1.3cm]}

\begin{tikzpicture}[node distance=8mm]
\node[draw=DeepBlue,very thick,rounded corners=3mm,fill=SoftBlue,
      minimum width=12.8cm,minimum height=3.2cm,align=center] (a)
{\Large\bfseries فصل اول\\[2mm]
\large تصویر به‌مثابه اندازه‌گیری، تابع و عملگر\\[1mm]
\normalsize از تشکیل فیزیکی تصویر تا صورت‌بندی مسائل معکوس};
\node[below=7mm of a,align=center,text width=13cm] {
این کتاب مبانی کلاسیک را حذف نمی‌کند؛ آن‌ها را با زبان ریاضی واحدی بازسازی می‌کند که
به پردازش آماری، بهینه‌سازی، یادگیری عمیق، مدل‌های مولد، مدل‌های بنیادی و تصویربرداری محاسباتی متصل می‌شود.};
\end{tikzpicture}

\vfill
{\large نسخه آغازین درس ـ تابستان ۱۴۰۵ / ۲۰۲۶}\\[0.4cm]
{\normalsize قابل کامپایل با \lr{TeX Live + TeXstudio + XeLaTeX}}\\[0.8cm]
{\color{DeepBlue}\rule{0.88\textwidth}{1.1pt}}
\end{titlepage}

\frontmatter

% -----------------------------------------------------------------------------
\chapter*{پیشگفتار}
\addcontentsline{toc}{chapter}{پیشگفتار}

پردازش تصویر در آستانه یک تغییر ساده از «روش‌های کلاسیک» به «یادگیری عمیق» نیست؛
بلکه در حال تبدیل‌شدن به دانشی است که در آن پنج لایه به‌طور هم‌زمان حضور دارند:
فیزیک تشکیل تصویر، ریاضیات نمایش و تبدیل، آمار و عدم‌قطعیت، بهینه‌سازی و یادگیری از داده،
و در نهایت مهندسی سامانه‌ای که باید در جهان واقعی قابل اعتماد، سریع و بازتولیدپذیر باشد.

در یک نگاه سطحی ممکن است تصور شود که با وجود مدل‌های بزرگ، دیگر نیازی به مطالعه دقیق
نمونه‌برداری، تبدیل فوریه، کانولوشن، هندسه تصویری، تخمین آماری یا مسائل معکوس نیست.
این تصور خطرناک است. مدل‌های جدید، قوانین فیزیکی حسگر، محدودیت اطلاعات، ناپایداری وارون‌سازی،
سوگیری داده و هزینه محاسبات را حذف نمی‌کنند؛ تنها ابزار تازه‌ای برای مدل‌کردن «پیشین» و
ساخت نمایش‌های عمومی‌تر فراهم می‌آورند. ظهور معماری‌های ترنسفورمری در بینایی،
یادگیری خودنظارتی در مقیاس بزرگ، مدل‌های قابل‌پرامپت برای قطعه‌بندی و مدل‌های دیفیوژن،
بیانگر آن است که مهندس آینده باید هم زبان پیکسل و عملگر را بداند و هم زبان توکن، نمایش،
امتیاز و پرامپت را \cite{dosovitskiy2021vit,he2022mae,kirillov2023sam,ho2020ddpm,oquab2024dinov2}.

از سوی دیگر، تصویربرداری محاسباتی مرز میان «دوربین» و «الگوریتم» را از بین برده است:
سامانه اپتیکی، الگوی نمونه‌برداری و الگوریتم بازسازی به‌صورت مشترک طراحی می‌شوند و تصویر
ممکن است مستقیماً روی حسگر تشکیل نشود، بلکه از اندازه‌گیری‌های کدگذاری‌شده بازسازی گردد
\cite{computationalimaging2026,zuo2022optical}. بنابراین، کتابی برای دهه آینده نباید از
فیلترگذاری آغاز کند؛ باید از این پرسش آغاز کند که «تصویر چیست و چگونه به داده تبدیل می‌شود؟»

این متن برای دانشجوی کارشناسی ارشد نوشته شده است، اما انتظار ریاضی آن عمداً بالاست.
هر فصل قرار است یک جزوه نسبتاً مستقل، یک مجموعه ویدئوی کوتاه، یک آزمایش محاسباتی،
یک تمرین اثباتی و یک مسئله مهندسی واقعی تولید کند. تجهیز هر دانشجو به
\lr{TeXstudio + TeX Live}، محیط برنامه‌نویسی پایتون، و یک دستیار کدنویسی و گفت‌وگویی
مانند \lr{Codex} و \lr{ChatGPT}، آزمایشگاه پردازش تصویر را از یک اتاق محدود دانشگاهی
به خانه دانشجو منتقل می‌کند؛ اما مسئولیت صحت ریاضی، طراحی آزمایش و تفسیر نتیجه همچنان
بر عهده پژوهشگر است.

\enlargethispage{3\baselineskip}
\begin{keybox}[فلسفه مرکزی کتاب]
در سراسر کتاب، هر روش با سه سؤال بازسازی می‌شود:
\begin{enumerate}
\item داده چگونه و با چه عملگری تولید شده است؟
\item چه اطلاعاتی از مجهول در اندازه‌گیری باقی مانده و چه اطلاعاتی از دست رفته است؟
\item کدام پیشین، قید، داده آموزشی یا مدل مولد می‌تواند فضای پاسخ‌ها را محدود کند؟
\end{enumerate}
\end{keybox}

% -----------------------------------------------------------------------------
\chapter*{چگونه از این کتاب استفاده کنیم؟}
\addcontentsline{toc}{chapter}{چگونه از این کتاب استفاده کنیم؟}

هر فصل از شش لایه تشکیل می‌شود:
\begin{enumerate}
\item \textbf{لایه شهودی:} صورت مسئله، مثال واقعی و تصویر ذهنی؛
\item \textbf{لایه ریاضی:} تعریف‌ها، قضایا، اثبات‌ها و تحلیل پایداری؛
\item \textbf{لایه الگوریتمی:} تبدیل مدل به یک روش عددی قابل اجرا؛
\item \textbf{لایه محاسباتی:} پیاده‌سازی کنترل‌شده و آزمون واحد؛
\item \textbf{لایه پژوهشی:} مقایسه روش‌های کلاسیک و جدید و آشکارکردن شکاف‌ها؛
\item \textbf{لایه مهندسی:} زمان اجرا، حافظه، خطا، بازتولیدپذیری و گزارش حرفه‌ای.
\end{enumerate}

\section*{قرارداد آموزشی پیشنهادی}
برای هر فصل، دانشجو باید چهار خروجی تحویل دهد:
\begin{itemize}
\item یک گزارش کوتاه \LaTeX{} که روابط اصلی را با بیان خود بازنویسی کند؛
\item یک مخزن کد با فایل \lr{README}، محیط قابل بازتولید و آزمون‌های پایه؛
\item یک آزمایش که در آن یک فرض مهم عمداً شکسته شود؛
\item یک ویدئوی پنج‌دقیقه‌ای که نتیجه را بدون پنهان‌کردن شکست‌ها توضیح دهد.
\end{itemize}

\section*{ابزارهای پایه}
\begin{longtable}{p{0.21\textwidth}p{0.7\textwidth}}
\toprule
\textbf{ابزار} & \textbf{کارکرد در درس}\\
\midrule
\endhead
\lr{TeX Live + TeXstudio} & نگارش روابط، گزارش آزمایش، ساخت شکل‌های برداری و تولید نسخه نهایی کتابچه هر فصل.\\
\lr{Python} & پیاده‌سازی عددی با \lr{NumPy, SciPy, scikit-image, OpenCV, Matplotlib} و در فصل‌های یادگیری با \lr{PyTorch}.\\
\lr{Git} & ثبت تغییرات، بازگشت به نسخه‌های سالم، مقایسه آزمایش‌ها و تحویل حرفه‌ای.\\
\lr{Codex/ChatGPT} & تولید اسکلت کد، نوشتن آزمون، توضیح خطا و پیشنهاد آزمایش؛ نه منبع نهایی حقیقت علمی.\\
\lr{NotebookLM} & گفت‌وگو با مجموعه محدود منابع انتخاب‌شده برای مرور یک زیرموضوع یا ساخت طرح ویدئو.\\
\bottomrule
\end{longtable}

% -----------------------------------------------------------------------------
\chapter*{نقشه کل کتاب و درس}
\addcontentsline{toc}{chapter}{نقشه کل کتاب و درس}

ساختار زیر عمداً از یک تقسیم‌بندی صرفاً تاریخی فاصله می‌گیرد. مبانی کلاسیک حفظ شده‌اند،
اما هر بخش به یک مفهوم ماندگار برای دهه آینده متصل است.

\begin{longtable}{p{0.06\textwidth}p{0.29\textwidth}p{0.55\textwidth}}
\toprule
\textbf{فصل} & \textbf{عنوان} & \textbf{هسته علمی و پل آینده}\\
\midrule
\endhead
۱ & تصویر به‌مثابه اندازه‌گیری، تابع و عملگر & تشکیل تصویر، فضاهای تابعی، عملگر پیشرو، بدوضعی و زبان واحد مسائل معکوس.\\
۲ & دیجیتالی‌شدن جهان دیداری & نمونه‌برداری، کوانتیزه‌سازی، علیاسینگ، دامنه دینامیکی، رنگ‌سنجی، \lr{HDR} و داده‌های چندطیفی.\\
۳ & سامانه‌های خطی و تحلیل فرکانسی & کانولوشن، پاسخ ضربه، فوریه، فیلترها، موجک و نمایش چندمقیاسی؛ مقدمه‌ای برای فیلترهای آموختنی.\\
۴ & تصویر تصادفی و مدل‌های نویز & میدان تصادفی، نویز گاوسی، پواسون، ضربه‌ای و وابسته به سیگنال؛ تخمین و آزمون آماری.\\
۵ & بهبود، حذف نویز و بازسازی & وینر، تیکانوف، تغییرات کلی، بهینه‌سازی پروگزیمال، \lr{PnP} و بازکردن تکرارها به شبکه.\\
۶ & هندسه تصویر و دوربین & هندسه تصویری، کالیبراسیون، هموگرافی، تبدیل مختصات، بازنمونه‌برداری و ثبت تصویر.\\
۷ & ساختارهای محلی و مورفولوژی & مشتق، لبه، گوشه، اسکلت، فاصله و مورفولوژی ریاضی؛ پیوند با سوگیری‌های ساختاری شبکه‌ها.\\
۸ & قطعه‌بندی کلاسیک و تغییراتی & آستانه‌گذاری، ناحیه، گراف، کانتور فعال، \lr{Mumford--Shah} و \lr{level set}.\\
۹ & ویژگی، تطبیق و بازشناسی کلاسیک & \lr{SIFT/HOG}, تطبیق مقاوم، \lr{RANSAC} و خط لوله کلاسیک قابل تفسیر.\\
۱۰ & شبکه‌های کانولوشنی و نمایش سلسله‌مراتبی & کانولوشن آموختنی، میدان پذیرش، نرمال‌سازی، باقیمانده، انتقال و کالیبراسیون.\\
۱۱ & توجه، ترنسفورمر و یادگیری خودنظارتی & توکن‌سازی تصویر، توجه، مدل‌سازی ماسک‌شده، یادگیری تقابلی و ویژگی‌های عمومی.\\
۱۲ & مدل‌های مولد و پیشین‌های آموخته‌شده & خودرمزگذار، \lr{VAE}, \lr{GAN}, مدل امتیاز و دیفیوژن؛ بازسازی و تولید کنترل‌شده.\\
۱۳ & مدل‌های بنیادی و بینایی قابل‌پرامپت & مدل‌های بینایی ـ زبان، قطعه‌بندی قابل‌پرامپت، سازگارسازی کم‌پارامتر و ارزیابی انتقال.\\
۱۴ & تصویربرداری محاسباتی و مدل‌های فیزیک‌آگاه & طراحی مشترک حسگر و الگوریتم، سنجش فشرده، فاز، توموگرافی، \lr{MRI} و یادگیری مقید به فیزیک.\\
۱۵ & ویدئو، حرکت، سه‌بعد و میدان‌های عصبی & جریان نوری، رهگیری، بازسازی سه‌بعدی، \lr{NeRF} و نمایش‌های ضمنی.\\
۱۶ & کیفیت، عدم‌قطعیت و اعتمادپذیری & معیارهای ادراکی، تغییر دامنه، خارج از توزیع، حمله، توهم بازسازی، انصاف و انسان در حلقه.\\
۱۷ & مهندسی حرفه‌ای سامانه‌های تصویری & داده، نسخه‌بندی، آزمون، پروفایل حافظه، کم‌سازی، استقرار لبه، گزارش مدل و بازتولیدپذیری.\\
۱۸ & استودیوی پایانی & پروژه کامل در پزشکی، صنعت، سنجش از دور، اسناد، رباتیک یا هنر با دفاع علمی و مهندسی.\\
\bottomrule
\end{longtable}

\begin{futurebox}
ترنسفورمرهای بینایی نشان دادند که تصویر می‌تواند به دنباله‌ای از وصله‌ها تبدیل شود؛
خودرمزگذارهای ماسک‌شده و \lr{DINOv2} نشان دادند که نمایش‌های قوی را می‌توان بدون برچسب مستقیم آموخت؛
\lr{SAM} مفهوم قطعه‌بندی قابل‌پرامپت را برجسته کرد؛ و مدل‌های دیفیوژن، حذف نویز را به هسته یک مدل مولد تبدیل کردند
\cite{dosovitskiy2021vit,he2022mae,oquab2024dinov2,kirillov2023sam,ho2020ddpm}.
این تحولات در فصل‌های بعد مطالعه می‌شوند، اما زبان ریاضی آن‌ها از همین فصل آغاز می‌شود.
\end{futurebox}

\section*{طرح پیشنهادی یک نیم‌سال ۱۴ هفته‌ای}
\begin{longtable}{p{0.09\textwidth}p{0.35\textwidth}p{0.45\textwidth}}
\toprule
\textbf{هفته} & \textbf{درس} & \textbf{خروجی کارگاهی}\\
\midrule
\endhead
۱ & فصل ۱: تصویر، اندازه‌گیری و عملگر & شبیه‌ساز ساده تشکیل تصویر و گزارش بدوضعی.\\
۲ & فصل ۲: نمونه‌برداری، کوانتیزه‌سازی و رنگ & آزمایش علیاسینگ و دامنه دینامیکی.\\
۳--۴ & فصل ۳: کانولوشن، فوریه و چندمقیاسی & طراحی فیلتر در فضا و فرکانس و تحلیل مرزی.\\
۵ & فصل ۴: نویز و آمار تصویر & تخمین پارامتر نویز از داده خام.\\
۶--۷ & فصل ۵: بازسازی و مسائل معکوس & مقایسه وینر، تیکانوف، \lr{TV} و یک پیشین آموخته‌شده.\\
۸ & فصل ۶ و ۷: هندسه و ساختار محلی & ثبت دو تصویر و تحلیل خطای درون‌یابی.\\
۹ & فصل ۸ و ۹: قطعه‌بندی و ویژگی کلاسیک & ساخت خط لوله کلاسیک قابل تفسیر.\\
۱۰ & فصل ۱۰: \lr{CNN} & آموزش یک مدل کوچک با تحلیل میدان پذیرش.\\
۱۱ & فصل ۱۱: ترنسفورمر و خودنظارتی & استخراج و نمایش ویژگی‌های مدل ازپیش‌آموزش‌دیده.\\
۱۲ & فصل ۱۲ و ۱۳: دیفیوژن و مدل بنیادی & آزمایش بازسازی/قطعه‌بندی با پرامپت و تحلیل شکست.\\
۱۳ & فصل ۱۴ و ۱۶: فیزیک، کیفیت و اعتماد & سنجش سازگاری داده و آشکارسازی توهم.\\
۱۴ & فصل ۱۷ و استودیوی پایانی & دفاع پروژه، گزارش بازتولیدپذیری و نمایش زنده.\\
\bottomrule
\end{longtable}

\tableofcontents

\mainmatter

% =============================================================================
\chapter{تصویر به‌مثابه اندازه‌گیری، تابع و عملگر}
\label{chap:image-as-measurement}

\begin{learningbox}
پس از مطالعه این فصل، دانشجو باید بتواند:
\begin{enumerate}
\item تصویر را پیش از دیجیتالی‌شدن به‌صورت یک میدان پیوسته فضایی ـ طیفی ـ زمانی تعریف کند؛
\item زنجیره تشکیل تصویر را از تابش صحنه تا عدد ذخیره‌شده در حافظه مدل کند؛
\item یک تصویر دیجیتال را به‌صورت تابع، ماتریس، بردار و تانسور نمایش دهد و میان این نمایش‌ها تبدیل انجام دهد؛
\item عملگرهای خطی، ناوردا نسبت به انتقال و غیرخطی را تشخیص دهد؛
\item مسئله پردازش تصویر را در قالب \(\vect y=\op A(\vect x)+\vect\eta\) بنویسد؛
\item خوش‌وضعی یا بدوضعی مسئله را با هسته، مقادیر منفرد و حساسیت به نویز تحلیل کند؛
\item نقش قید، منظم‌ساز، پیشین احتمالاتی و مدل آموخته‌شده را در بازسازی توضیح دهد؛
\item یک آزمایش عددی بازتولیدپذیر برای تاری، نویز و وارون‌سازی طراحی کند.
\end{enumerate}
\end{learningbox}

\section{چرا فصل نخست باید از «اندازه‌گیری» آغاز شود؟}

در بسیاری از کتاب‌های کلاسیک، نخستین تعریف تصویر چنین است:
«تصویر تابعی دوبعدی مانند \(f(x,y)\) است و مقدار تابع شدت روشنایی را نشان می‌دهد.»
این تعریف برای آغاز مناسب است، اما برای مهندس دهه آینده کافی نیست؛ زیرا:
\begin{itemize}
\item مقدار ثبت‌شده الزاماً شدت واقعی صحنه نیست؛ حاصل اپتیک، پاسخ طیفی حسگر، نوردهی، نویز و پردازش داخلی دوربین است؛
\item داده ممکن است دوبعدی نباشد: تصویر رنگی، ویدئو، داده حجمی، چندطیفی، قطبشی یا رویدادمحور تانسوری با ابعاد بیشتر است؛
\item در بسیاری از سامانه‌ها، حسگر اصلاً تصویر مطلوب را مستقیماً اندازه نمی‌گیرد؛ برای مثال در توموگرافی، \lr{MRI}، بازیابی فاز یا تصویربرداری فشرده، داده مجموعه‌ای از اندازه‌گیری‌های غیرمستقیم است؛
\item الگوریتم مدرن گاهی نه شدت پیکسل، بلکه توکن، ویژگی، توزیع، امتیاز یا نمایش ضمنی را پردازش می‌کند.
\end{itemize}

بنابراین نقطه شروع مناسب‌تر این است:

\begin{keybox}[تعریف راهبردی]
یک سامانه تصویری، جهان یا شیء ناشناخته \(x\) را به‌وسیله یک زنجیره فیزیکی و محاسباتی
به داده مشاهده‌شده \(y\) تبدیل می‌کند. پردازش تصویر علم مدل‌کردن، تحلیل، تغییر، تفسیر
یا وارون‌سازی این نگاشت است.
\end{keybox}

این تعریف، هم افزایش کنتراست ساده و هم بازسازی \lr{MRI}، هم قطعه‌بندی و هم تولید تصویر
با دیفیوژن را در یک چارچوب قرار می‌دهد.

\section{صحنه و تصویر پیوسته}

\subsection{تابش صحنه}

فرض کنید نقطه‌ای از جهان با مختصات فضایی \((X,Y,Z)\)، در طول موج \(\lambda\) و زمان \(t\)
تابشی تولید یا بازتاب می‌کند. یک مدل غنی از صحنه می‌تواند به‌صورت میدان تابش
\begin{equation}
L=L(X,Y,Z,\theta,\phi,\lambda,t)
\end{equation}
نوشته شود که علاوه بر مکان، جهت انتشار \((\theta,\phi)\)، طول موج و زمان را نیز در بر دارد.
دوربین معمولی تنها فرافکنی بسیار محدودی از این میدان پُربعد را ثبت می‌کند.

برای یک سطح لامبرتی ساده، می‌توان وابستگی جهت را تقریباً حذف کرد و در یک مدل ابتدایی نوشت
\begin{equation}
L(X,Y,Z,\lambda,t)=E(X,Y,Z,\lambda,t)\,\rho(X,Y,Z,\lambda),
\end{equation}
که در آن \(E\) روشنایی تابیده‌شده و \(\rho\) بازتابندگی سطح است. حتی در این مدل ساده نیز
یک پیکسل صرفاً «رنگ شیء» نیست؛ حاصل ضرب نور و ویژگی سطح است.

\begin{warningbox}
برابردانستن مقدار پیکسل با یک ویژگی ذاتی جسم، منشأ بسیاری از شکست‌های الگوریتمی است.
تغییر نور، نوردهی، پاسخ دوربین و پردازش خودکار می‌تواند مقدار پیکسل را بدون تغییر جسم عوض کند.
\end{warningbox}

\subsection{تابع تصویر تک‌رنگ}

پس از فرافکنی هندسی و اثر اپتیک، یک تصویر ایده‌آل تک‌رنگ را می‌توان به‌صورت
\begin{equation}
f:\Omega\subset\R^2\longrightarrow \R_{\ge 0},
\qquad (x,y)\longmapsto f(x,y)
\end{equation}
تعریف کرد. دامنه \(\Omega\) ناحیه حسگر یا صفحه تصویر و \(f(x,y)\) انرژی یا شدت نوری مؤثر است.

اما این تابع غالباً یک میانگین طیفی است. اگر \(s(\lambda)\) حساسیت طیفی حسگر باشد،
\begin{equation}
f(x,y)=\int_{\Lambda} L_{\mathrm{image}}(x,y,\lambda)\,s(\lambda)\,d\lambda.
\label{eq:spectral-integration}
\end{equation}
در دوربین رنگی سه پاسخ متفاوت به‌کار می‌رود:
\begin{equation}
f_c(x,y)=\int_{\Lambda}L_{\mathrm{image}}(x,y,\lambda)s_c(\lambda)d\lambda,
\qquad c\in\{R,G,B\}.
\end{equation}
پس تصویر رنگی تابعی برداری است:
\begin{equation}
f:\Omega\to\R^3.
\end{equation}

\subsection{تصویر، ویدئو و داده چندوجهی}

با افزودن زمان، ویدئو به‌صورت
\begin{equation}
f:\Omega\times[0,T]\to\R^C
\end{equation}
تعریف می‌شود. داده حجمی پزشکی ممکن است
\begin{equation}
f:\Omega_x\times\Omega_y\times\Omega_z\to\R
\end{equation}
باشد و تصویر فراطیفی به شکل
\begin{equation}
f:\Omega_x\times\Omega_y\times\Lambda\to\R
\end{equation}
مدل شود. در نتیجه «تصویر» در این کتاب نامی برای یک آرایه دوبعدی ثابت نیست؛
نمایشی نمونه‌برداری‌شده از یک میدان یا شیء پُربعد است.

\section{زنجیره تشکیل تصویر}

\subsection{مدل لایه‌ای}

شکل \ref{fig:image-formation-chain} زنجیره‌ای را نشان می‌دهد که از صحنه آغاز و به داده دیجیتال ختم می‌شود.

\begin{figure}[H]
\centering
\begin{tikzpicture}[
box/.style={draw=DeepBlue,thick,rounded corners,fill=SoftBlue,minimum height=12mm,
             text width=25mm,align=center},
arr/.style={-{Latex[length=2.5mm]},thick,DeepBlue},
lab/.style={font=\small,align=center,text width=24mm}]
\node[box] (scene) {صحنه و تابش\\\(x\)};
\node[box,right=7mm of scene] (geom) {هندسه و حرکت\\\(\op P\)};
\node[box,right=7mm of geom] (optics) {اپتیک و تاری\\\(\op H\)};
\node[box,right=7mm of optics] (sensor) {انتگرال‌گیری حسگر\\\(\op S\)};
\node[box,right=7mm of sensor] (adc) {نویز و تبدیل دیجیتال\\\(\op Q\)};
\node[box,right=7mm of adc,fill=SoftGold] (data) {داده ذخیره‌شده\\\(y\)};
\draw[arr] (scene)--(geom);
\draw[arr] (geom)--(optics);
\draw[arr] (optics)--(sensor);
\draw[arr] (sensor)--(adc);
\draw[arr] (adc)--(data);
\node[below=7mm of sensor,draw=Accent,dashed,rounded corners,text width=9.5cm,align=center]
{مدل پیشرو: \(y=\op Q\!\left(\op S\!\left(\op H\!\left(\op P(x)\right)\right)+\eta\right)\)};
\end{tikzpicture}
\caption{زنجیره مفهومی تشکیل تصویر. هر بلوک می‌تواند خطی، غیرخطی، قطعی یا تصادفی باشد.}
\label{fig:image-formation-chain}
\end{figure}

در ساده‌ترین حالت تک‌رنگ، بدون اعوجاج هندسی و با اپتیک خطی ناوردا نسبت به انتقال،
تصویر پیوسته روی حسگر برابر است با
\begin{equation}
g(x,y)=(h*f)(x,y)+\eta(x,y),
\label{eq:continuous-blur}
\end{equation}
که در آن \(h\) تابع پخش نقطه‌ای یا \lr{PSF} و \(*\) کانولوشن دوبعدی است:
\begin{equation}
(h*f)(x,y)=\int_{\R^2}h(\xi,\zeta)f(x-\xi,y-\zeta)d\xi d\zeta.
\end{equation}

\subsection{تابع پخش نقطه‌ای}

اگر ورودی اپتیکی یک نقطه ایده‌آل در مبدأ باشد، خروجی سامانه \(h(x,y)\) است.
برای یک سامانه خطی، هر تصویر را می‌توان جمع پیوسته‌ای از نقاط دانست؛ در نتیجه پاسخ کل
جمع پاسخ به این نقاط خواهد بود و کانولوشن ظاهر می‌شود.

\begin{definition}[سامانه خطی]
عملگر \(\op H\) خطی است هرگاه برای هر \(f_1,f_2\) و ضرایب \(a,b\) داشته باشیم
\begin{equation}
\op H(af_1+bf_2)=a\op H(f_1)+b\op H(f_2).
\end{equation}
\end{definition}

\begin{definition}[ناوردایی نسبت به انتقال]
اگر \(T_{\Delta}f(x,y)=f(x-\Delta_x,y-\Delta_y)\)، عملگر \(\op H\) نسبت به انتقال ناورداست هرگاه
\begin{equation}
\op H(T_{\Delta}f)=T_{\Delta}\op H(f).
\end{equation}
\end{definition}

\begin{theorem}[نمایش کانولوشنی]
هر سامانه پیوسته خطی و ناوردا نسبت به انتقال، تحت شرایط منظم‌بودن مناسب، با کانولوشن ورودی و پاسخ ضربه نمایش داده می‌شود:
\begin{equation}
\op H(f)=h*f,
\qquad h=\op H(\delta).
\end{equation}
\end{theorem}

\begin{proof}
با نمایش تابع به‌صورت برهم‌نهی ضربه‌های انتقال‌یافته
\begin{equation}
f(x,y)=\int_{\R^2}f(\xi,\zeta)\delta(x-\xi,y-\zeta)d\xi d\zeta
\end{equation}
و استفاده از خطی‌بودن، عملگر را به داخل انتگرال می‌بریم. ناوردایی نسبت به انتقال می‌دهد
\begin{equation}
\op H\{\delta(x-\xi,y-\zeta)\}=h(x-\xi,y-\zeta).
\end{equation}
پس
\begin{equation}
\op H(f)=\int f(\xi,\zeta)h(x-\xi,y-\zeta)d\xi d\zeta=h*f.
\end{equation}
\end{proof}

\subsection{انتگرال‌گیری روی سطح پیکسل}

پیکسل یک نقطه ریاضی نیست. پیکسل واقعی در یک ناحیه \(P_{mn}\) و بازه نوردهی
\([t_0,t_0+T_e]\) فوتون جمع می‌کند. مقدار ایده‌آل پیش از نویز می‌تواند به‌صورت
\begin{equation}
\mu_{mn}=\int_{t_0}^{t_0+T_e}\int_{P_{mn}}g(x,y,t)\,a_{mn}(x,y)\,dxdy\,dt
\label{eq:pixel-integration}
\end{equation}
نوشته شود که \(a_{mn}\) پاسخ مکانی پیکسل است.
این رابطه روشن می‌کند چرا حرکت در زمان نوردهی تاری می‌سازد و چرا کوچک‌کردن پیکسل
به‌تنهایی همیشه کیفیت را افزایش نمی‌دهد: فوتون کمتر، نسبت سیگنال به نویز ضعیف‌تر و حساسیت بیشتر به پراش ممکن است رخ دهد.

\subsection{نویز فوتونی و خوانش}

در مدل فوتونی، تعداد فوتون‌های ثبت‌شده اغلب با متغیر پواسون تقریب زده می‌شود:
\begin{equation}
Z_{mn}\sim\operatorname{Poisson}(\mu_{mn}).
\end{equation}
خوانش الکترونیکی نیز ممکن است نویز گاوسی بیفزاید:
\begin{equation}
R_{mn}\sim\mathcal N(0,\sigma_r^2),
\qquad Y_{mn}^{\mathrm{analog}}=g_s Z_{mn}+R_{mn}+b,
\end{equation}
که \(g_s\) بهره حسگر و \(b\) بایاس است. پس نویز همیشه جمعی، گاوسی و مستقل از سیگنال نیست.

\section{از تابع پیوسته تا تصویر دیجیتال}

\subsection{نمونه‌برداری مکانی}

اگر شبکه نمونه‌برداری با فواصل \(\Delta_x,\Delta_y\) باشد، نمونه‌ها به شکل
\begin{equation}
g[m,n]=g(m\Delta_x,n\Delta_y),
\qquad m=0,\dots,M-1,\quad n=0,\dots,N-1
\end{equation}
تعریف می‌شوند. مدل دقیق‌تر، طبق رابطه \eqref{eq:pixel-integration}، میانگین یا انتگرال روی سطح پیکسل است.

تابع نمونه‌برداری شانه‌ای دوبعدی را تعریف می‌کنیم:
\begin{equation}
\operatorname{III}_{\Delta_x,\Delta_y}(x,y)
=\sum_{m\in\mathbb Z}\sum_{n\in\mathbb Z}
\delta(x-m\Delta_x)\delta(y-n\Delta_y).
\end{equation}
نمونه‌برداری ایده‌آل برابر است با
\begin{equation}
g_s(x,y)=g(x,y)\operatorname{III}_{\Delta_x,\Delta_y}(x,y).
\end{equation}
در فصل دوم نشان خواهیم داد که ضرب در حوزه مکان، به تکرار طیف در حوزه فرکانس منجر می‌شود
و هم‌پوشانی این نسخه‌ها علیاسینگ را ایجاد می‌کند.

\subsection{کوانتیزه‌سازی}

پس از نمونه‌برداری، دامنه پیوسته یا پُردقت به مجموعه‌ای محدود از سطوح نگاشت می‌شود.
برای کوانتیزه‌ساز یکنواخت \(B\)-بیتی با بازه \([a,b]\)، اندازه گام تقریباً
\begin{equation}
\Delta_q=\frac{b-a}{2^B-1}
\end{equation}
است و
\begin{equation}
Y[m,n]=Q_B(g[m,n]+\eta[m,n]).
\end{equation}
اگر اشباع رخ ندهد و خطای کوانتیزه‌سازی به‌طور تقریبی یکنواخت باشد،
\begin{equation}
e_q\sim\mathcal U\left(-\frac{\Delta_q}{2},\frac{\Delta_q}{2}\right),
\qquad \operatorname{Var}(e_q)=\frac{\Delta_q^2}{12}.
\end{equation}
این تقریب در نزدیکی سطوح ثابت، بیت کم یا پردازش غیرخطی ممکن است معتبر نباشد.

\section{چهار نمایش ریاضی یک تصویر دیجیتال}

فرض کنید تصویر خاکستری \(M\times N\) باشد.

\subsection{تصویر به‌صورت تابع روی شبکه}
\begin{equation}
f:\{0,\dots,M-1\}\times\{0,\dots,N-1\}\to\R.
\end{equation}
این نمایش برای نوشتن فیلترهای محلی و هندسه پیکسل مناسب است.

\subsection{تصویر به‌صورت ماتریس}
\begin{equation}
\vect F=
\begin{bmatrix}
f[0,0]&\cdots&f[0,N-1]\\
\vdots&\ddots&\vdots\\
f[M-1,0]&\cdots&f[M-1,N-1]
\end{bmatrix}\in\R^{M\times N}.
\end{equation}
در این نمایش، رتبه، مقادیر منفرد، نرم هسته‌ای و ساختار بلوکی قابل مطالعه‌اند.

\subsection{تصویر به‌صورت بردار}

با عملگر \(\operatorname{vec}\)، ستون‌ها یا سطرهای تصویر پشت سر هم قرار می‌گیرند:
\begin{equation}
\vect f=\operatorname{vec}(\vect F)\in\R^{MN}.
\end{equation}
این نمایش برای نوشتن عملگرهای خطی به‌صورت ماتریس سودمند است:
\begin{equation}
\vect y=\vect A\vect f+\vect\eta.
\end{equation}
اما ماتریس \(\vect A\) غالباً نباید به‌طور صریح ساخته شود؛ برای تصویر مگاپیکسلی، ابعاد آن می‌تواند نجومی باشد.

\subsection{تصویر به‌صورت تانسور}

یک دسته \(B\)-تایی از تصاویر رنگی با ابعاد \(H\times W\) در یادگیری عمیق معمولاً به شکل
\begin{equation}
\vect X\in\R^{B\times C\times H\times W}
\end{equation}
ذخیره می‌شود. ترتیب محورها یک قرارداد نرم‌افزاری است، نه حقیقت فیزیکی.
برای ویدئو، محور زمان افزوده می‌شود و برای داده فراطیفی، محور طیف معنای فیزیکی دارد.

\begin{warningbox}
سه خطای بسیار رایج در کدهای دانشجویی عبارت‌اند از: جابه‌جایی ارتفاع و عرض، جابه‌جایی کانال و دسته،
و استفاده ناخواسته از نوع داده صحیح هنگام انجام محاسبات اعشاری. هر سه می‌توانند خروجی ظاهراً قابل قبول اما علمی نادرست بسازند.
\end{warningbox}

\section{تصویر در فضاهای برداری و تابعی}

\subsection{ضرب داخلی و انرژی}

در فضای گسسته \(\R^{M\times N}\)، ضرب داخلی فرابنیوس برابر است با
\begin{equation}
\inner{\vect F}{\vect G}_F
=\sum_{m=0}^{M-1}\sum_{n=0}^{N-1}F_{mn}G_{mn}
=\operatorname{trace}(\vect F\trans\vect G).
\end{equation}
نرم متناظر
\begin{equation}
\norm{\vect F}_F=\sqrt{\inner{\vect F}{\vect F}_F}
\end{equation}
انرژی اقلیدسی تصویر را اندازه می‌گیرد.

در حالت پیوسته، فضای \(L^2(\Omega)\) شامل توابعی است که
\begin{equation}
\norm{f}_{L^2}^2=\int_\Omega\abs{f(x,y)}^2dxdy<\infty.
\end{equation}

\subsection{نرم‌های \(\ell_p\)}

برای بردار \(\vect f\in\R^d\):
\begin{equation}
\norm{\vect f}_p=\left(\sum_{i=1}^d\abs{f_i}^p\right)^{1/p},\qquad p\ge1,
\end{equation}
و
\begin{equation}
\norm{\vect f}_\infty=\max_i\abs{f_i}.
\end{equation}
\(\ell_2\) انرژی خطا را برجسته می‌کند، \(\ell_1\) نسبت به خطاهای بزرگ مقاوم‌تر است و
همچنین می‌تواند تنکی را تشویق کند، و \(\ell_\infty\) بدترین خطا را کنترل می‌کند.

\begin{example}[دو خطا با انرژی برابر]
دو بردار خطا را در نظر بگیرید:
\begin{equation}
\vect e_1=(1,1,1,1),\qquad \vect e_2=(2,0,0,0).
\end{equation}
هر دو نرم \(\ell_2\) برابر ۲ دارند، اما
\begin{equation}
\norm{\vect e_1}_1=4,\qquad \norm{\vect e_2}_1=2,
\end{equation}
و
\begin{equation}
\norm{\vect e_1}_\infty=1,\qquad \norm{\vect e_2}_\infty=2.
\end{equation}
پس واژه «خطا» بدون مشخص‌کردن معیار ناقص است.
\end{example}

\subsection{فضاهای هموار و تغییرات کلی}

تصاویر طبیعی در لبه‌ها ناپیوستگی دارند، بنابراین فرض همواری مربعی سراسری همیشه مناسب نیست.
نیم‌نرم مرتبه اول را می‌توان به شکل
\begin{equation}
\norm{f}_{H^1}^2=\norm{f}_{L^2}^2+\norm{\nabla f}_{L^2}^2
\end{equation}
نوشت. این معیار گرادیان‌های بزرگ را به‌صورت مربعی جریمه می‌کند و ممکن است لبه‌ها را بیش از حد صاف کند.

تغییرات کلی همسانگرد در حالت پیوسته برابر است با
\begin{equation}
\operatorname{TV}(f)=\int_\Omega\norm{\nabla f(x,y)}_2dxdy.
\end{equation}
نسخه گسسته رایج:
\begin{equation}
\operatorname{TV}(\vect F)=\sum_{m,n}
\sqrt{(D_x\vect F)_{mn}^2+(D_y\vect F)_{mn}^2+\varepsilon^2}.
\end{equation}
جریمه خطی اندازه گرادیان نسبت به جریمه مربعی، لبه‌های بزرگ را کمتر سرکوب می‌کند.

\section{عملگرها: واحد سازنده پردازش تصویر}

\subsection{عملگر خطی و الحاقی}

برای فضاهای ضرب داخلی \(\mathcal X,\mathcal Y\)، عملگر خطی
\(\op A:\mathcal X\to\mathcal Y\) دارای عملگر الحاقی \(\op A^*\) است اگر
\begin{equation}
\inner{\op A x}{y}_{\mathcal Y}=\inner{x}{\op A^*y}_{\mathcal X}
\end{equation}
برای همه \(x,y\) برقرار باشد. در حالت ماتریسی حقیقی، الحاقی همان ترانهاده است.

عملگر الحاقی در گرادیان تابع داده نقش بنیادی دارد. اگر
\begin{equation}
\mathcal D(x)=\frac12\norm{\op A x-y}_2^2,
\end{equation}
آنگاه
\begin{equation}
\nabla\mathcal D(x)=\op A^*(\op A x-y).
\label{eq:data-gradient}
\end{equation}
این رابطه در کمترین مربعات، بازفرافکنی توموگرافی، شبکه‌های بازشده و بسیاری از روش‌های مدرن تکرار می‌شود.

\subsection{چند عملگر مهم}

\begin{longtable}{p{0.2\textwidth}p{0.34\textwidth}p{0.36\textwidth}}
\toprule
\textbf{عملگر} & \textbf{مدل} & \textbf{نمونه کاربرد}\\
\midrule
\endhead
همانی & \(\op A(x)=x\) & حذف نویز؛ مجهول مستقیماً با نویز مشاهده می‌شود.\\
تاری & \(\op A(x)=h*x\) & تاری خارج از فوکوس، حرکت یا پراش.\\
زیرنمونه‌برداری & \(\op A(x)=\op D x\) & ابرتفکیک و داده کم‌وضوح.\\
پوشاندن & \(\op A(x)=\vect M\odot x\) & پرکردن نواحی گمشده و پیکسل‌های خراب.\\
فوریه جزئی & \(\op A(x)=\vect P\vect F x\) & بازسازی \lr{MRI}.\\
رادون & \(\op A(x)=\mathcal R x\) & توموگرافی و \lr{CT}.\\
قدر مطلق فوریه & \(\op A(x)=\abs{\vect F x}^2\) & بازیابی فاز؛ عملگر غیرخطی.\\
شبکه دوربین & \(\op A(x;\theta)\) & مدل پاسخ دوربین، دموُزاییک و پردازش داخلی.\\
\bottomrule
\end{longtable}

\subsection{عملگرهای غیرخطی}

گاما، اشباع، آستانه‌گذاری، میانه، قدر مطلق، کوانتیزه‌سازی و شبکه عصبی نمونه‌هایی از عملگر غیرخطی‌اند.
برای مثال تبدیل گاما
\begin{equation}
y=x^{1/\gamma}
\end{equation}
جمع‌پذیر نیست. فیلتر میانه نیز به ترتیب آماری همسایگی وابسته است و نمی‌توان آن را با یک هسته ثابت کانولوشن نمایش داد.

\section{مدل عمومی پیشرو}

مدل بنیادی این کتاب چنین است:
\begin{equation}
\vect y=\op A(\vect x;\vect\theta)+\vect\eta,
\label{eq:general-forward}
\end{equation}
که در آن:
\begin{itemize}
\item \(\vect x\) کمیت مطلوب، مانند تصویر پاک، تصویر پُروضوح، حجم یا پارامتر صحنه است؛
\item \(\vect y\) داده اندازه‌گیری‌شده است؛
\item \(\op A\) مدل پیشرو فیزیکی یا محاسباتی است؛
\item \(\vect\theta\) پارامترهای کالیبراسیون یا مجهول عملگر است؛
\item \(\vect\eta\) نویز، خطای مدل و اغتشاش‌های ثبت‌نشده را دربرمی‌گیرد.
\end{itemize}

گاهی خود \(\op A\) نیز نامعلوم است. در تاری کور، هم تصویر \(x\) و هم هسته \(h\) باید تخمین زده شوند:
\begin{equation}
y=h*x+\eta.
\end{equation}
این مسئله بسیار دشوارتر از دکانولوشن با هسته معلوم است، زیرا زوج‌های متعدد \((x,h)\) می‌توانند خروجی مشابه بسازند.

\begin{keybox}[اصل سازگاری داده]
هر خروجی بازسازی‌شده \(\widehat x\) باید پس از عبور از مدل پیشرو، اندازه‌گیری واقعی را توضیح دهد:
\begin{equation}
\op A(\widehat x)\approx y.
\end{equation}
کیفیت ظاهری به‌تنهایی کافی نیست. تصویری زیبا که با داده فیزیکی سازگار نیست، ممکن است «توهم بازسازی» باشد.
\end{keybox}

\section{مسئله معکوس و مفهوم بدوضعی}

\subsection{شرایط هادامار}

\begin{definition}[مسئله خوش‌وضع]
یک مسئله به معنای هادامار خوش‌وضع است اگر:
\begin{enumerate}
\item پاسخ وجود داشته باشد؛
\item پاسخ یکتا باشد؛
\item پاسخ به داده وابستگی پیوسته داشته باشد.
\end{enumerate}
اگر هر کدام نقض شود، مسئله بدوضع است.
\end{definition}

در پردازش تصویر، شرط سوم بسیار مهم است. ممکن است وارون ریاضی وجود داشته باشد، اما نویز بسیار کوچک در داده به خطای بسیار بزرگ در پاسخ تبدیل شود.

\subsection{تحلیل مقادیر منفرد}

برای مدل خطی گسسته
\begin{equation}
\vect y=\vect A\vect x+\vect\eta,
\end{equation}
تجزیه مقدار منفرد را بنویسید:
\begin{equation}
\vect A=\vect U\vect\Sigma\vect V\trans,
\qquad
\vect\Sigma=\operatorname{diag}(\sigma_1,\ldots,\sigma_r),
\end{equation}
که \(\sigma_1\ge\cdots\ge\sigma_r>0\) هستند. وارون شبه‌مور برای داده بدون نویز
\begin{equation}
\widehat{\vect x}_{\mathrm{pinv}}
=\vect A^\dagger\vect y
=\sum_{i=1}^{r}\frac{\vect u_i\trans\vect y}{\sigma_i}\vect v_i
\end{equation}
است. با جای‌گذاری \(\vect y=\vect A\vect x+\vect\eta\):
\begin{equation}
\widehat{\vect x}_{\mathrm{pinv}}
=\vect V\vect V\trans\vect x
+\sum_{i=1}^{r}\frac{\vect u_i\trans\vect\eta}{\sigma_i}\vect v_i.
\end{equation}
اگر \(\sigma_i\) کوچک باشد، مؤلفه نویز بر \(\sigma_i\) تقسیم و بزرگ می‌شود.

\begin{proposition}[کران حساسیت]
اگر \(\vect A\) وارون‌پذیر باشد، تغییر پاسخ نسبت به اغتشاش داده از رابطه
\begin{equation}
\norm{\delta\vect x}_2
\le \norm{\vect A^{-1}}_2\norm{\delta\vect y}_2
=\frac{1}{\sigma_{\min}(\vect A)}\norm{\delta\vect y}_2
\end{equation}
پیروی می‌کند.
\end{proposition}

عدد وضعیت
\begin{equation}
\kappa_2(\vect A)=\frac{\sigma_{\max}(\vect A)}{\sigma_{\min}(\vect A)}
\end{equation}
مقیاس تقویت نسبی خطا را بیان می‌کند. در دکانولوشن، فرکانس‌هایی که پاسخ انتقال اپتیک در آن‌ها کوچک است، همان جهت‌های کم‌مقدار منفردند.

\subsection{عدم یکتایی و هسته}

اگر بردار ناصفر \(\vect z\in\ker(\vect A)\) باشد، آنگاه
\begin{equation}
\vect A(\vect x+\vect z)=\vect A\vect x.
\end{equation}
پس اندازه‌گیری نمی‌تواند میان \(\vect x\) و \(\vect x+\vect z\) تمایز بگذارد.
هیچ الگوریتمی نمی‌تواند اطلاعاتی را که عملگر به‌طور کامل حذف کرده است، تنها از داده بازیابی کند؛
باید قید یا دانش پیشین افزوده شود.

\section{از برازش داده تا منظم‌سازی}

\subsection{کمترین مربعات}

برای نویز گاوسی سفید، یک تخمین طبیعی عبارت است از
\begin{equation}
\widehat{\vect x}_{\mathrm{LS}}
\in\arg\min_{\vect x}\frac12\norm{\vect A\vect x-\vect y}_2^2.
\end{equation}
شرط مرتبه اول:
\begin{equation}
\vect A\trans(\vect A\widehat{\vect x}-\vect y)=0.
\end{equation}
اگر \(\vect A\trans\vect A\) وارون‌پذیر باشد،
\begin{equation}
\widehat{\vect x}=(\vect A\trans\vect A)^{-1}\vect A\trans\vect y.
\end{equation}
اما تشکیل معادلات نرمال عدد وضعیت را تقریباً مربع می‌کند و برای مسائل بزرگ مناسب نیست.

\subsection{منظم‌سازی تیکانوف}

برای کنترل ناپایداری، مسئله را به شکل
\begin{equation}
\widehat{\vect x}_\lambda
=\arg\min_{\vect x}
\left\{
\frac12\norm{\vect A\vect x-\vect y}_2^2
+\frac{\lambda}{2}\norm{\vect L\vect x}_2^2
\right\}
\label{eq:tikhonov}
\end{equation}
می‌نویسیم. پاسخ، در صورت وارون‌پذیری، برابر است با
\begin{equation}
\widehat{\vect x}_\lambda
=(\vect A\trans\vect A+\lambda\vect L\trans\vect L)^{-1}\vect A\trans\vect y.
\end{equation}

برای \(\vect L=\vect I\)، فیلتر روی مؤلفه‌های منفرد چنین است:
\begin{equation}
\widehat{\vect x}_\lambda
=\sum_i
\underbrace{\frac{\sigma_i}{\sigma_i^2+\lambda}}_{\text{عامل بازسازی}}
(\vect u_i\trans\vect y)\vect v_i.
\end{equation}
در وارون مستقیم عامل \(1/\sigma_i\) بود، اما اکنون برای \(\sigma_i\ll\sqrt\lambda\)،
تقویت محدود می‌شود. بهای این پایداری، بایاس است.

\subsection{منظم‌سازی تغییرات کلی}

برای حفظ لبه‌ها می‌توان نوشت
\begin{equation}
\widehat{\vect x}
=\arg\min_{\vect x}
\frac12\norm{\vect A\vect x-\vect y}_2^2
+\lambda\operatorname{TV}(\vect x).
\label{eq:tv}
\end{equation}
این مسئله معمولاً پاسخ بسته ندارد و با روش‌های پروگزیمال، \lr{ADMM} یا الگوریتم‌های اولیه ـ دوگان حل می‌شود.

\subsection{صورت احتمالاتی و تخمین MAP}

اگر درست‌نمایی داده \(p(\vect y\mid\vect x)\) و پیشین \(p(\vect x)\) معلوم باشد،
قاعده بیز می‌دهد
\begin{equation}
p(\vect x\mid\vect y)
\propto p(\vect y\mid\vect x)p(\vect x).
\end{equation}
تخمین بیشینه پسین:
\begin{align}
\widehat{\vect x}_{\mathrm{MAP}}
&=\arg\max_{\vect x}p(\vect x\mid\vect y)\\
&=\arg\min_{\vect x}
\left[-\log p(\vect y\mid\vect x)-\log p(\vect x)\right].
\end{align}
اگر نویز گاوسی باشد، جمله نخست به کمترین مربعات تبدیل می‌شود. اگر پیشین
\begin{equation}
p(\vect x)\propto\exp(-\lambda\mathcal R(\vect x))
\end{equation}
باشد، منفی لگاریتم پیشین همان منظم‌ساز است. پس منظم‌سازی و استنباط بیزی دو زبان مرتبط‌اند.

\section{پیشین‌های دست‌ساز، آموخته‌شده و مولد}

\subsection{پیشین دست‌ساز}

نمونه‌ها عبارت‌اند از همواری، تنکی گرادیان، تنکی موجک، خودشباهتی غیرمحلی، رتبه کم و محدودیت بازه شدت.
مزیت آن‌ها تفسیرپذیری و امکان تحلیل ریاضی است؛ ضعف آن‌ها محدودیت در مدل‌کردن ساختار پیچیده تصویر است.

\subsection{پیشین ضمنی در حذف‌کننده نویز}

چارچوب \lr{Plug-and-Play} یک حل‌کننده سازگاری داده را با یک حذف‌کننده نویز ترکیب می‌کند،
حتی اگر حذف‌کننده صریحاً گرادیان یک منظم‌ساز شناخته‌شده نباشد \cite{venkatakrishnan2013pnp}.
به‌صورت شماتیک:
\begin{align}
\vect z^{k+1}&\approx\arg\min_z\mathcal D(z;y)+\frac{\rho}{2}\norm{z-\vect v^k}_2^2,\\
\vect x^{k+1}&=D_\sigma(\vect z^{k+1}),
\end{align}
که \(D_\sigma\) حذف‌کننده نویز است.

\subsection{پیشین مولد و مدل امتیاز}

مدل مولد می‌کوشد توزیع تصاویر را بیاموزد. در مدل‌های امتیاز، کمیت
\begin{equation}
s_\theta(\vect x,t)\approx\nabla_{\vect x}\log p_t(\vect x)
\end{equation}
یاد گرفته می‌شود. مدل‌های دیفیوژن با افزودن تدریجی نویز و یادگیری فرایند معکوس،
حذف نویز را به سازوکار تولید تبدیل می‌کنند \cite{ho2020ddpm}.
در بازسازی، چنین مدلی می‌تواند یک پیشین بسیار قوی باشد؛ اما اگر سازگاری داده ضعیف شود،
جزئیات محتمل ولی نادرست تولید می‌شوند.

\begin{futurebox}
در دهه آینده، مرز میان «فیلتر»، «منظم‌ساز»، «حذف‌کننده نویز» و «مدل مولد» کم‌رنگ‌تر می‌شود.
با این حال، معادله \eqref{eq:general-forward} همچنان معیار اتصال مدل آموخته‌شده به واقعیت فیزیکی خواهد بود.
مدل قوی‌تر، نیاز به مدل پیشرو و سنجش عدم‌قطعیت را از بین نمی‌برد؛ اهمیت آن‌ها را بیشتر می‌کند.
\end{futurebox}

\section{از پیکسل تا توکن و نمایش}

ترنسفورمر بینایی تصویر \(H\times W\times C\) را به وصله‌های \(P\times P\) تقسیم می‌کند.
تعداد توکن‌ها
\begin{equation}
N=\frac{HW}{P^2}
\end{equation}
است. هر وصله برداری و با نگاشت خطی به یک بردار \(d\)-بعدی تبدیل می‌شود:
\begin{equation}
\vect z_i=\vect E\operatorname{vec}(\vect X_i)+\vect p_i,
\end{equation}
که \(\vect p_i\) کد موقعیت است. توجه خودی برای ماتریس توکن‌ها \(\vect Z\) به شکل
\begin{equation}
\operatorname{Attention}(\vect Q,\vect K,\vect V)
=\operatorname{softmax}\left(\frac{\vect Q\vect K\trans}{\sqrt{d_k}}\right)\vect V
\end{equation}
تعریف می‌شود.

این نمایش جایگزین حقیقت فیزیکی تصویر نیست. توکن‌سازی خود نوعی نمونه‌برداری و فشرده‌سازی است.
اندازه وصله، کد موقعیت و معماری توجه تعیین می‌کنند چه جزئیاتی حفظ یا حذف شود.
مدل‌های خودنظارتی مانند \lr{MAE} بخشی از وصله‌ها را پنهان و بازسازی می‌کنند و مدل‌هایی مانند
\lr{DINOv2} نمایش‌های عمومی را بدون برچسب مستقیم می‌آموزند \cite{he2022mae,oquab2024dinov2}.

\section{سه سطح مسئله: پردازش، تحلیل و فهم}

\begin{enumerate}
\item \textbf{پردازش سطح پایین:} ورودی و خروجی تصویرند؛ مانند حذف نویز، رفع تاری، ابرتفکیک و اصلاح رنگ.
\item \textbf{تحلیل سطح میانی:} خروجی ساختار است؛ مانند لبه، ناحیه، عمق، حرکت، تطبیق یا ماسک.
\item \textbf{فهم سطح بالا:} خروجی مفهوم، تصمیم یا پاسخ زبانی است؛ مانند تشخیص، شرح تصویر و استدلال دیداری.
\end{enumerate}

مدل‌های بنیادی این مرزها را مخلوط کرده‌اند. یک نمایش واحد ممکن است برای دسته‌بندی، عمق و قطعه‌بندی استفاده شود؛
یک مدل قابل‌پرامپت می‌تواند بدون آموزش اختصاصی ماسک تولید کند؛ و مدل‌های مولد می‌توانند هم تصویر بسازند و هم
به‌عنوان پیشین بازسازی به کار روند. با وجود این، معیار موفقیت در هر سطح متفاوت است و نباید کیفیت ادراکی،
دقت پارامتری و صحت معنایی با یکدیگر خلط شوند.

\section{مثال یکپارچه: رفع تاری حرکت}

فرض کنید تصویر پاک \(\vect x\) با هسته حرکت \(\vect h\) تار و سپس با نویز آلوده شود:
\begin{equation}
\vect y=\vect H\vect x+\vect\eta.
\end{equation}

\subsection{راه‌حل خام}
در حوزه فوریه، اگر \(Y=HX+N\)، وارون مستقیم
\begin{equation}
\widehat X=\frac{Y}{H}
\end{equation}
است. اگر \(H(\omega)\) در برخی فرکانس‌ها نزدیک صفر باشد، نویز شدیداً تقویت می‌شود.

\subsection{تیکانوف در حوزه فرکانس}
برای \(\vect L=\vect I\):
\begin{equation}
\widehat X_\lambda(\omega)
=\frac{\overline{H(\omega)}}{\abs{H(\omega)}^2+\lambda}Y(\omega).
\end{equation}
این فیلتر در فرکانس‌های مطمئن تقریباً وارون و در فرکانس‌های ضعیف سرکوب‌کننده است.

\subsection{پیشین لبه}
در روش \lr{TV}، انتظار داریم تصویر در نواحی صاف گرادیان کم و در مرزها تعداد محدودی گرادیان بزرگ داشته باشد.
حل غیرخطی‌تر اما لبه‌حافظ‌تر است.

\subsection{پیشین آموخته‌شده}
یک شبکه بازسازی یا مدل دیفیوژن می‌تواند ساختارهای پیچیده‌تر را بازسازی کند؛ ولی باید بررسی شود آیا جزئیات
بازسازی‌شده واقعاً توسط داده پشتیبانی می‌شوند یا از توزیع آموزشی «حدس» زده شده‌اند.

\section{آزمایش محوری فصل}

\begin{workshopbox}[از عملگر پیشرو تا بازسازی پایدار]
\textbf{هدف:} دانشجو باید یک تصویر آزمایشی بسازد، آن را با سه عملگر تخریب کند و نشان دهد چرا وارون مستقیم شکست می‌خورد.

\textbf{مراحل:}
\begin{enumerate}
\item یک تصویر خاکستری استاندارد و یک تصویر مصنوعی شامل ناحیه صاف، لبه تیز و بافت بسازید.
\item سه \lr{PSF} تعریف کنید: گاوسی، حرکت خطی و تاری خارج از فوکوس.
\item برای هر کدام نویز گاوسی و پواسون را جداگانه بیفزایید.
\item پاسخ فرکانسی هسته و توزیع مقادیر منفرد عملگر را رسم کنید.
\item وارون مستقیم، تیکانوف و \lr{TV} را اجرا کنید.
\item پارامتر منظم‌سازی را تنها با نگاه‌کردن به تصویر انتخاب نکنید؛ منحنی خطای بازسازی، باقیمانده داده و حساسیت را گزارش کنید.
\item آزمایش را با هسته اشتباه تکرار کنید تا اثر خطای مدل مشخص شود.
\item برای هر خروجی، هم معیار پیکسلی و هم سازگاری داده \(\norm{\op A(\widehat x)-y}\) را ثبت کنید.
\end{enumerate}
\end{workshopbox}

\subsection{اسکلت الگوریتم تیکانوف با گرادیان مزدوج}

به‌جای ساخت ماتریس بزرگ، عملگر پیشرو و الحاقی را به‌صورت تابع پیاده‌سازی کنید.
برای حل
\begin{equation}
(\op A^*\op A+\lambda\op L^*\op L)x=\op A^*y
\end{equation}
از گرادیان مزدوج استفاده کنید.

\begin{LTR}
\begin{lstlisting}[caption={Matrix-free Tikhonov solver skeleton}]
import numpy as np
from scipy.sparse.linalg import LinearOperator, cg


def solve_tikhonov(y, A, AT, L, LT, shape, lam, rtol=1e-6):
    """Solve (A^T A + lam L^T L)x = A^T y matrix-free."""
    n = int(np.prod(shape))

    def matvec(x_vec):
        x = x_vec.reshape(shape)
        out = AT(A(x)) + lam * LT(L(x))
        return np.asarray(out, dtype=np.float64).ravel()

    normal_op = LinearOperator((n, n), matvec=matvec, dtype=np.float64)
    rhs = np.asarray(AT(y), dtype=np.float64).ravel()
    x_vec, info = cg(normal_op, rhs, rtol=rtol)
    if info != 0:
        raise RuntimeError(f"CG did not converge; info={info}")
    return x_vec.reshape(shape)
\end{lstlisting}
\end{LTR}

\subsection{آزمون‌هایی که باید نوشته شوند}
\begin{itemize}
\item \textbf{آزمون الحاقی:}
\begin{equation}
\frac{\abs{\inner{\op A x}{y}-\inner{x}{\op A^*y}}}
{\max(1,\abs{\inner{\op A x}{y}},\abs{\inner{x}{\op A^*y}})}<10^{-6}.
\end{equation}
\item \textbf{آزمون حد بدون تخریب:} اگر \(A=I\) و نویز صفر باشد، بازسازی باید ورودی را با خطای عددی بازگرداند.
\item \textbf{آزمون یکنواختی:} کانولوشن هسته نرمال‌شده با تصویر ثابت باید تصویر ثابت را، دور از مرزها، حفظ کند.
\item \textbf{آزمون بازتولیدپذیری:} بذر تصادفی، نسخه کتابخانه‌ها و پارامترها ذخیره شوند.
\end{itemize}

\begin{keybox}[پرامپت پیشنهادی برای دستیار کدنویسی]
یک پروژه پایتون بازتولیدپذیر برای آزمایش فصل اول بساز. عملگرهای تاری گاوسی، حرکت و خارج از فوکوس را
به‌صورت توابع \lr{forward} و \lr{adjoint} پیاده‌سازی کن و آزمون الحاقی بنویس. نویز گاوسی و پواسون را
با بذر ثابت اضافه کن. وارون مستقیم فرکانسی، تیکانوف ماتریس‌آزاد با گرادیان مزدوج و \lr{TV} را پیاده‌سازی کن.
برای هر حالت، \lr{MSE, PSNR, SSIM}، خطای سازگاری داده، زمان و حافظه را ذخیره کن. آزمایش خطای مدل را با
استفاده از هسته بازسازی متفاوت از هسته تولید داده اجرا کن. تمام شکل‌ها و جدول‌ها باید از فایل نتایج تولید شوند؛
هیچ عددی در گزارش به‌صورت دستی نوشته نشود. فایل‌های \lr{requirements.txt}، \lr{README.md} و آزمون‌های
\lr{pytest} نیز ایجاد شوند.
\end{keybox}

\section{چگونه یک آزمایش را علمی گزارش کنیم؟}

یک شکل زیبا به‌تنهایی آزمایش نیست. گزارش باید دست‌کم شامل موارد زیر باشد:
\begin{enumerate}
\item مدل دقیق تولید داده و واحد پارامترها؛
\item داده واقعی یا مصنوعی و دلیل انتخاب آن؛
\item محدوده پارامترها و روش انتخاب ابرپارامتر؛
\item معیارهای کیفیت، پارامتری، سازگاری داده و هزینه محاسباتی؛
\item میانگین و انحراف معیار روی چند تکرار در حضور تصادف؛
\item تحلیل شکست، نه فقط بهترین حالت؛
\item کد، بذر، محیط اجرا و مسیر تولید خودکار جدول‌ها.
\end{enumerate}

\begin{warningbox}
انتخاب بهترین \(\lambda\) با استفاده از تصویر حقیقت و سپس گزارش همان خطا به‌عنوان نتیجه نهایی،
اگر حقیقت در کاربرد واقعی در دسترس نیست، ارزیابی خوش‌بینانه است. برای داده مصنوعی باید نقش آن را
«کران تشخیصی» دانست؛ برای کاربرد واقعی باید روش انتخاب مستقلی مانند اعتبارسنجی، اصل ناسازگاری،
\lr{SURE} یا معیار مبتنی بر فیزیک طراحی شود.
\end{warningbox}

\section{تمرین‌های ریاضی}

\subsection*{تمرین ۱: خطی‌بودن و ناوردایی}
برای هر عملگر زیر تعیین کنید خطی و/یا ناوردا نسبت به انتقال است:
\begin{enumerate}
\item \(\op A(f)=f^2\)،
\item \(\op A(f)=h*f\)،
\item \(\op A(f)(x,y)=x f(x,y)\)،
\item فیلتر میانه \(3\times3\)،
\item \(\op A(f)=f-\operatorname{mean}(f)\).
\end{enumerate}
برای هر پاسخ اثبات یا مثال نقض ارائه کنید.

\subsection*{تمرین ۲: الحاقی زیرنمونه‌برداری}
فرض کنید \(D\) فقط پیکسل‌های زوج را نگه می‌دارد. عملگر الحاقی \(D^*\) را استخراج کنید و نشان دهید
\(D^*\) در مکان‌های مشاهده‌نشده صفر درج می‌کند؛ سپس تفاوت آن را با درون‌یابی توضیح دهید.

\subsection*{تمرین ۳: فیلتر تیکانوف}
با استفاده از \lr{SVD} نشان دهید حل \eqref{eq:tikhonov} برای \(L=I\) دارای عامل فیلتر
\begin{equation}
\phi_i(\lambda)=\frac{\sigma_i^2}{\sigma_i^2+\lambda}
\end{equation}
در بازسازی مؤلفه واقعی است. حدهای \(\lambda\to0\) و \(\lambda\to\infty\) را تفسیر کنید.

\subsection*{تمرین ۴: MAP با نویز پواسون}
اگر \(y_i\sim\operatorname{Poisson}((Ax)_i)\)، منفی لگاریتم درست‌نمایی را تا ثابت‌های مستقل از \(x\)
به‌دست آورید و نشان دهید چرا کمترین مربعات معمولی دقیقاً متناظر با این نویز نیست.

\subsection*{تمرین ۵: ناپایداری}
یک ماتریس قطری \(A=\operatorname{diag}(1,10^{-1},10^{-2},10^{-4})\) بسازید. برای یک بردار معلوم،
نویز با اندازه \(10^{-4}\) اضافه کنید و وارون مستقیم را تحلیل کنید. جهت خطرناک نویز کدام است؟

\subsection*{تمرین ۶: هسته عملگر}
برای عملگر مشتق گسسته مرتبه اول، هسته را تعیین کنید. توضیح دهید چرا بازسازی یک سیگنال از مشتق آن
تا یک ثابت نامعین است و چه قیدی یکتایی را بازمی‌گرداند.

\section{تمرین‌های پژوهشی و مهندسی}

\begin{enumerate}
\item یک دوربین گوشی را در حالت \lr{RAW} و \lr{JPEG} روی صحنه ثابت آزمایش کنید. هیستوگرام، اشباع،
شارپ‌سازی و تغییرات کانال رنگ را مقایسه کنید و توضیح دهید کدام بلوک‌های شکل \ref{fig:image-formation-chain}
در خروجی \lr{JPEG} پنهان شده‌اند.
\item برای یک تصویر پزشکی یا صنعتی، یک بازسازی «ظاهراً بهتر اما ناسازگار با داده» بسازید و نشان دهید
چرا معیار ادراکی به‌تنهایی خطرناک است.
\item اندازه حافظه لازم برای ساخت صریح ماتریس تاری یک تصویر \(1024\times1024\) را محاسبه کنید.
سپس پیاده‌سازی ماتریس‌آزاد را با آن مقایسه کنید.
\item یک حذف‌کننده نویز ازپیش‌آموزش‌دیده را در چارچوب \lr{PnP} قرار دهید و رفتار آن را هنگام تغییر دامنه تصویر،
نوع نویز و شدت تخریب بررسی کنید.
\item یک مدل ویژگی عمومی مانند \lr{DINOv2} را روی تصویر و نسخه‌های تار، نویزی و تغییررنگ‌یافته آن اجرا کنید.
فاصله ویژگی‌ها را با فاصله پیکسلی مقایسه و درباره ناوردایی و حساسیت بحث کنید.
\end{enumerate}

\section{تقسیم فصل به ویدئوهای کوتاه}

\begin{videobox}
\begin{longtable}{p{0.08\textwidth}p{0.36\textwidth}p{0.45\textwidth}}
\toprule
\textbf{بخش} & \textbf{موضوع} & \textbf{خروجی روی تخته یا رایانه}\\
\midrule
\endhead
۱ & چرا «تصویر = آرایه پیکسل» کافی نیست؟ & معرفی مدل پیشرو و سه مثال غیرمستقیم.\\
۲ & تابش، بازتاب و پاسخ طیفی & استخراج رابطه انتگرال طیفی یک کانال.\\
۳ & اپتیک و \lr{PSF} & ساخت شهودی کانولوشن از پاسخ نقطه.\\
۴ & پیکسل واقعی، نوردهی و نویز & مدل پواسون ـ گاوسی حسگر.\\
۵ & نمونه‌برداری و کوانتیزه‌سازی & آزمایش بصری کاهش نرخ نمونه و بیت.\\
۶ & تابع، ماتریس، بردار و تانسور & تبدیل \lr{shape/reshape/vec} با آزمون.\\
۷ & فضاها و نرم‌ها & مقایسه \(\ell_1,\ell_2,\ell_\infty,TV\).\\
۸ & عملگر، الحاقی و آزمون الحاقی & پیاده‌سازی \lr{forward/adjoint}.\\
۹ & مدل عمومی مسئله معکوس & نوشتن مدل برای تاری، ماسک، \lr{MRI} و فاز.\\
۱۰ & بدوضعی با \lr{SVD} & نمایش تقویت نویز در مقدار منفرد کوچک.\\
۱۱ & تیکانوف و سازش بایاس ـ واریانس & استخراج پاسخ بسته و عامل فیلتر.\\
۱۲ & \lr{TV} و لبه & تفاوت جریمه گرادیان مربعی و خطی.\\
۱۳ & بیز، پیشین و \lr{MAP} & تبدیل منفی لگاریتم به تابع هدف.\\
۱۴ & پیشین آموخته‌شده، \lr{PnP} و دیفیوژن & توضیح ظرفیت و خطر توهم.\\
۱۵ & از پیکسل تا توکن & وصله‌سازی و معادله توجه.\\
۱۶ & طراحی آزمایش نهایی & معیارها، خطای مدل و بازتولیدپذیری.\\
\bottomrule
\end{longtable}
\end{videobox}

\section{خلاصه فصل}

\begin{enumerate}
\item تصویر مشاهده‌شده نسخه مستقیم و بی‌واسطه جهان نیست؛ خروجی یک زنجیره اندازه‌گیری است.
\item مدل پیشرو \(y=\op A(x)+\eta\) ستون فقرات تحلیل پردازش تصویر است.
\item نمایش‌های تابعی، ماتریسی، برداری و تانسوری مکمل‌اند و هر کدام ابزارهای خاصی فراهم می‌کنند.
\item کانولوشن نتیجه خطی‌بودن و ناوردایی نسبت به انتقال است، نه صرفاً یک دستور نرم‌افزاری.
\item بدوضعی از نبود پاسخ، نبود یکتایی یا ناپایداری ناشی می‌شود؛ مقدار منفرد کوچک نویز را تقویت می‌کند.
\item منظم‌سازی اطلاعاتی فراتر از داده وارد می‌کند و میان برازش و پایداری سازش می‌سازد.
\item پیشین می‌تواند دست‌ساز، احتمالاتی، ضمنی در حذف‌کننده نویز یا آموخته‌شده توسط مدل مولد باشد.
\item در مدل‌های مدرن، پیکسل ممکن است به توکن و نمایش تبدیل شود، اما محدودیت اطلاعات و فیزیک همچنان پابرجاست.
\item کیفیت ادراکی بدون سازگاری داده، به‌ویژه در کاربردهای حساس، کافی نیست.
\item مهندس حرفه‌ای باید مدل، الگوریتم، کد، آزمایش و گزارش را به یک زنجیره بازتولیدپذیر تبدیل کند.
\end{enumerate}

\section{واژگان کلیدی فصل}

\begin{tabularx}{\textwidth}{>{\bfseries}p{0.28\textwidth}X}
مدل پیشرو & نگاشت مجهول فیزیکی یا تصویری به داده اندازه‌گیری‌شده.\\
مسئله معکوس & تخمین علت یا مجهول از روی اندازه‌گیری‌های مستقیم یا غیرمستقیم.\\
تابع پخش نقطه‌ای & پاسخ سامانه تصویری به یک منبع نقطه‌ای؛ هسته کانولوشن در سامانه \lr{LSI}.\\
الحاقی & عملگری که ضرب داخلی خروجی و داده را به ضرب داخلی ورودی و بازفرافکنی تبدیل می‌کند.\\
بدوضعی & نقض وجود، یکتایی یا پایداری پاسخ.\\
منظم‌سازی & افزودن قید یا پیشین برای محدودکردن پاسخ و افزایش پایداری.\\
سازگاری داده & میزان انطباق اندازه‌گیری پیش‌بینی‌شده از پاسخ با داده واقعی.\\
توهم بازسازی & تولید جزئیات محتمل یا زیبا که به اندازه کافی توسط اندازه‌گیری پشتیبانی نمی‌شوند.\\
نمایش خودنظارتی & ویژگی آموخته‌شده بدون نیاز به برچسب مستقیم برای هر نمونه.\\
مدل بنیادی & مدل بزرگ ازپیش‌آموزش‌دیده که قابلیت انتقال به دامنه‌ها یا وظایف متعدد دارد.\\
\end{tabularx}

% -----------------------------------------------------------------------------
\backmatter
\renewcommand{\bibname}{منابع منتخب برای نقشه کتاب و فصل اول}
\addcontentsline{toc}{chapter}{منابع منتخب}
\begin{thebibliography}{99}
\begin{LTR}

\bibitem{gonzalez2018dip}
R. C. Gonzalez and R. E. Woods,
\emph{Digital Image Processing}, 4th ed., Pearson, 2018.

\bibitem{szeliski2022cv}
R. Szeliski,
\emph{Computer Vision: Algorithms and Applications}, 2nd ed., Springer, 2022.

\bibitem{bertero1998inverse}
M. Bertero and P. Boccacci,
\emph{Introduction to Inverse Problems in Imaging}, Institute of Physics Publishing, 1998.

\bibitem{hansen2010inverse}
P. C. Hansen,
\emph{Discrete Inverse Problems: Insight and Algorithms}, SIAM, 2010.

\bibitem{beck2017optimization}
A. Beck,
\emph{First-Order Methods in Optimization}, SIAM, 2017.

\bibitem{rudin1992tv}
L. I. Rudin, S. Osher, and E. Fatemi,
``Nonlinear total variation based noise removal algorithms,''
\emph{Physica D}, vol. 60, pp. 259--268, 1992.

\bibitem{venkatakrishnan2013pnp}
S. V. Venkatakrishnan, C. A. Bouman, and B. Wohlberg,
``Plug-and-Play Priors for Model Based Reconstruction,''
\emph{IEEE Global Conference on Signal and Information Processing}, pp. 945--948, 2013.

\bibitem{dosovitskiy2021vit}
A. Dosovitskiy et al.,
``An Image is Worth 16x16 Words: Transformers for Image Recognition at Scale,''
\emph{International Conference on Learning Representations}, 2021.

\bibitem{he2022mae}
K. He, X. Chen, S. Xie, Y. Li, P. Doll\'ar, and R. Girshick,
``Masked Autoencoders Are Scalable Vision Learners,''
\emph{CVPR}, pp. 16000--16009, 2022.

\bibitem{zamir2022restormer}
S. W. Zamir et al.,
``Restormer: Efficient Transformer for High-Resolution Image Restoration,''
\emph{CVPR}, pp. 5728--5739, 2022.

\bibitem{ho2020ddpm}
J. Ho, A. Jain, and P. Abbeel,
``Denoising Diffusion Probabilistic Models,''
\emph{Advances in Neural Information Processing Systems}, vol. 33, 2020.

\bibitem{kirillov2023sam}
A. Kirillov et al.,
``Segment Anything,''
\emph{ICCV}, pp. 4015--4026, 2023.

\bibitem{oquab2024dinov2}
M. Oquab et al.,
``DINOv2: Learning Robust Visual Features without Supervision,''
\emph{Transactions on Machine Learning Research}, 2024.

\bibitem{zuo2022optical}
C. Zuo et al.,
``Deep learning in optical metrology: a review,''
\emph{Light: Science \& Applications}, vol. 11, article 39, 2022.

\bibitem{computationalimaging2026}
Nature Portfolio,
``Computational Imaging,''
\emph{Nature Index Topic Overview}, accessed July 2026.

\bibitem{zhou2024segfm}
T. Zhou et al.,
``Image Segmentation in Foundation Model Era: A Survey,''
arXiv:2408.12957, 2024.

\end{LTR}
\end{thebibliography}

\end{document}
